% Source: https://artofproblemsolving.com/community/c3h568592_mathcounts_national_2009_sprint
% Stable identity: aops-topic-568592-q23
\textbf{Claim:} $D$ is the centroid. If $[BCD]=\dfrac{1}{3}[ABC]$, then the distance from $D$ to $BC$ is one third of the distance from $A$ to $BC$, so $D$ lies on the line parallel to $BC$ one third of the way from $BC$ to $A$; similarly for the other two sides. The centroid $G$ satisfies these conditions, because it lies on each median two thirds of the way from the vertex, and only one point lies on two of these parallel lines. Hence $D=G$:
\[
(m,n)=\left(\dfrac{5+3+6}{3},\dfrac{8-2+1}{3}\right)=\left(\dfrac{14}{3},\dfrac{7}{3}\right).
\]

\begin{asy}
usepackage("amsmath");
size(6.5cm);
defaultpen(fontsize(10pt));
pair A=(5,8),B=(3,-2),C=(6,1);
pair G=(A+B+C)/3;
draw(A--B--C--cycle,linewidth(0.9));
draw(A--(B+C)/2,red+dashed); draw(B--(A+C)/2,red+dashed); draw(C--(A+B)/2,red+dashed);
dot(A);dot(B);dot(C);dot(G);
label("$A(5,8)$",A,N); label("$B(3,-2)$",B,S); label("$C(6,1)$",C,E);
label("$D$",G,dir(12));
label("\mbox{}",(min(currentpicture,user=true)+max(currentpicture,user=true))/2,black); // keeps the text after the figure black when LaTeX typesets the labels inline
\end{asy}

Then $10m+n=\dfrac{140}{3}+\dfrac{7}{3}=\dfrac{147}{3}=\boxed{49}$.
